\chapter{光学}
\section{光線の幾何と場の位相}
幾何光学は光の波長より十分大きい構造で光線と像を扱い、波動光学は重ね合わせと位相を扱う。
レンズの実像・虚像は光線の交点の問題であり、回折・干渉の明暗は場の位相差の問題である。
同じ屈折率が両方へ関係するが、条件を取り違えないことが重要である。
光電効果と電子線回折の比較は現代物理の第\ref{theory:modern:duality}節へ分類する。
幾何から波動へという構成は大学の光学講義にも共通する\cite{mit_optics2014}。

\section{薄レンズの符号で実像と虚像を分ける}\label{theory:optics:lens}
近軸・薄レンズ近似において、凸レンズの焦点距離を \(f>0\)、入射側の実物体までを \(a>0\) とする。
出射側の像を \(b>0\)、入射側の虚像を \(b<0\) と定めれば
\[
 \frac1f=\frac1a+\frac1b,\qquad b=\frac{fa}{a-f},\qquad M=-\frac ba
\]
である。\(M\) は像高／物体高であり、正なら正立、負なら倒立である。
焦点外 \(a>f\) は倒立実像、焦点内 \(0<a<f\) は正立虚像となる。
\(a=f\) なら出射光線は平行となり、有限の像面はない\cite{os_lens}。
「虚像だから小さい」という推論は正しくない。焦点内の凸レンズでは \(|M|>1\) となる。
実像の移動方向は
\[
 \frac{db}{da}=-\frac{f^2}{(a-f)^2}<0
\]
から直接判定する。焦点外の物体がレンズへ近付くと像は遠ざかる。
焦点を越えて物体を動かす場合、有限の実像がそのままレンズを通過して虚像になるわけではなく、無限遠を経由する。

\section{二枚のレンズは中間像の位置を保存する}\label{theory:optics:cascade}
第一レンズを \(x=0\)、第二を \(x=\ell\) とし、第一の像位置を \(x=b_1\) とする。
第二へ来る光は、第一の像位置から発散又はその位置へ収束する束である。
第二からの物体距離は
\[
 a_2=\ell-b_1
\]
であり、\(b_1<0\) の第一の虚像なら \(a_2>\ell>0\) となる。
第一の「虚像」が第二に対しても必ず「虚物体」というわけではない。
逆に、第一の実像予定位置が第二のさらに後方にあれば \(a_2<0\) の虚物体となる。
各レンズの式を順に解き、最後に同じ座標軸で位置を照合する。
レンズ間隔を各物体距離へそのまま置く方法では中間像の位置を失ってしまう。

\section{屈折率・全反射・色収差}\label{theory:optics:refraction}
屈折率は光の位相速度を真空の光速 \(c\) で割った逆数 \(n=c/v_{\mathrm{ph}}\) である。
境界に平行な波数成分の保存から \(n_1\sin\theta_1=n_2\sin\theta_2\) を得る。
高屈折率側から低屈折率側へ入射し、\(\sin\theta_1>n_2/n_1\) となると、
伝播する屈折波を作れず全反射となる。
したがって光ファイバのコアはクラッドより高い屈折率を持つ\cite{os_fiber}。
境界外にはエバネッセントな場があり、「外側の電磁場が必ず厳密に零」とはしない。
通常の透明域の正常分散では紫の屈折率が赤より大きい。
同形の薄い凸レンズを空気中に置くと、曲率に依存する定数 \(C>0\) に対して \(1/f=(n-1)C\) となるので、
紫は焦点距離が短く、赤は遠くへ集光する。これは正常分散の範囲であり、共鳴近傍の異常分散まで無条件に含めない。

\section{反射位相を経路差へ足す}\label{theory:optics:reflection}
損失のない誘電体へほぼ垂直に入射する電場の反射係数は \(r=(n_1-n_2)/(n_1+n_2)\) である。
低屈折率側から高屈折率側なら \(r<0\) で位相が \(\pi\) 反転し、逆なら反転しない。
膜の両面からの反射光を比較するときは、光学的経路差と反射位相を別に加える\cite{os_thinfilm}。
角度・偏光・吸収まで一般化すると反射係数はこれより複雑になる。
問題に角度がない場合の通常の説明を、全入射角に同じ位相反転が成立する普遍則として教えない。

\section{ニュートンリングの幾何と次数}\label{theory:optics:newton}
曲率半径 \(R\) の凸面を平面へ置くと、接点から半径 \(x\) における空気層の厚さは
\[
 t(x)=R-\sqrt{R^2-x^2}\simeq\frac{x^2}{2R}\quad(x\ll R)
\]
となる。垂直入射の反射光を観察し、一方だけ位相反転すると、暗環は \(2t=n\lambda\)、明環は
\[
 2t=(n+1/2)\lambda,\qquad x_n^2=(n+1/2)R\lambda,\qquad n=0,1,2,\ldots
\]
である。接点は暗く、最初の明環は \(x^2=R\lambda/2\) となる。
ここで \(n\) は零始まりの干渉次数である。
「第 \(m\) 番目」を \(m=1,2,\ldots\) とする序数なら \(n=m-1\) であり、\(x_m^2=(m-1/2)R\lambda\) と書く。
開始値を指定せず両方を同じ記号で使わない。
令和3年度の原典は「\(m\) 番目」とする一方、公式肢は零始まりの形であるため解釈疑義を残している。
同じ曲率・同じ次数で波長が長いと半径が増すことは、この開始値の違いによらない。

\section{回折格子：位置・幅・強度を分ける}\label{theory:optics:grating}
垂直入射、間隔 \(d\) の等間隔スリットから遠方へ向かう光の隣接経路差は \(d\sin\theta\) である。
全てのスリットが同位相となる主極大は
\[
 d\sin\theta=m\lambda,\qquad m=0,\pm1,\pm2,\ldots
\]
であり、スクリーン距離 \(L\) と小角近似で主極大間隔は \(\Delta x\simeq L\lambda/d\) となる\cite{os_grating}。
入射強度は線形領域で全体の明るさを変え、この位置を変えない。
スリット数 \(N\) の効果は、位相差 \(\delta=2\pi d\sin\theta/\lambda\) を用いて和を取ると分かる。
\[
 \left|\sum_{j=0}^{N-1}e^{ij\delta}\right|^2
 =\left(\frac{\sin(N\delta/2)}{\sin(\delta/2)}\right)^2
\]
であり、\(\delta=2\pi m\) は極限値 \(N^2\) の主極大である。
主極大間には \(N-1\) 個の零点と \(N-2\) 個の副極大がある。
本数を増すと主極大は狭くなり、同じ主極大間の細かい暗線は増える。
有限幅スリットの単スリット包絡も別に掛かる\cite{os_double}。
したがって「縞間隔」を主極大間隔と読むか、全ての暗線の間隔と読むかを区別する。
平成27年度の「明線（もしくは暗線）」という表現には、この適用範囲の疑義を保存した。
格子本数は \(1/d\) である。波長を nm、密度を本/mmで求める場合は、まず \(d\) を mmへ換算する。
第一次数の650 nmと30度なら \(d=1.3\,\mu\mathrm m\) で、約769本/mmとなる。
これは主極大の方位を測る問題であり、縞の幅や副極大を数える問題ではない。

\section{電磁波領域と光子エネルギー}\label{theory:optics:spectrum}
真空中では \(\lambda f=c\) である。
波長の長い順は、ここで問う4領域では赤外・可視・紫外・γである\cite{os_spectrum}。
光子エネルギー \(E=hf=hc/\lambda\) の大きい順なら逆になる。
ビーム全体のエネルギーは光子数にも依存するので、波長の短い光が任意の照射条件で必ず大きな総エネルギーを持つとは言えない。
この区別は分光、放射線、光電効果の問題へつながる。

\section{出題の型と原典確認の範囲}
像の問題は物体と像の位置の符号、干渉は反射位相・経路差・次数、格子は間隔と本数の逆数、
スペクトルは比較する量を先に確定する。
S7で光学の11登録問を確認し、電子線回折の1問を現代物理へ移した。
対象14問の解説初稿を本章へ置く。表の空欄や次数の疑義は解説と未解決表で追跡する。

\section{登録問題と解説初稿}
\subsection{平成23年度 問7}
\label{q:h23-reg-applied-07}

\noindent\textbf{原典：} sources/H23/17\_applied\_science.pdf、PDF物理3頁。

\textbf{公式正答：} 2（single）。

\par\textit{原典本文の校正と独自解説は未作成。分類のみの登録項目。}

\subsection{平成24年度 問3}
\label{q:h24-reg-applied-03}

\noindent\textbf{原典：} sources/H24/17\_applied\_science.pdf、PDF物理1頁。

\textbf{公式正答：} 1（single）。

\par\textit{原典本文の校正と独自解説は未作成。分類のみの登録項目。}

\subsection{平成24年度 問7}
\label{q:h24-reg-applied-07}

\noindent\textbf{原典：} sources/H24/17\_applied\_science.pdf、PDF物理2頁。

\textbf{公式正答：} 2（single）。

\par\textit{原典本文の校正と独自解説は未作成。分類のみの登録項目。}

\subsection{平成26年度 問7}
\label{q:h26-reg-applied-07}

\noindent\textbf{原典：} sources/H26/17\_applied\_science.pdf、PDF物理3頁。

\textbf{公式正答：} 2（single）。

\par\textit{原典本文の校正と独自解説は未作成。分類のみの登録項目。}

\subsection{平成27年度 問3}
\label{q:h27-reg-applied-03}

\noindent\textbf{原典：} sources/H27/17\_applied\_science.pdf、PDF物理1頁。

\textbf{公式正答：} 1（single）。

\par\textit{原典本文の校正と独自解説は未作成。分類のみの登録項目。}

\subsection{平成28年度 問5}
\label{q:h28-reg-applied-05}

\noindent\textbf{原典：} sources/H28/17\_applied\_science.pdf、PDF物理2頁。

\textbf{公式正答：} 3（single）。

\par\textit{原典本文の校正と独自解説は未作成。分類のみの登録項目。}

\subsection{令和元年度再試験 問3}
\label{q:r01-ret-applied-03}

\noindent\textbf{原典：} sources/R01\_retest/17\_applied\_science.pdf、PDF物理2頁。

\textbf{公式正答：} 5（single）。

\par\textit{原典本文の校正と独自解説は未作成。分類のみの登録項目。}

\subsection{令和2年度 問5}
\label{q:r02-reg-applied-05}

\noindent\textbf{原典：} sources/R02/17\_applied\_science.pdf、PDF物理2頁。

\textbf{公式正答：} 5（single）。

\par\textit{原典本文の校正と独自解説は未作成。分類のみの登録項目。}

\subsection{令和2年度 問7}
\label{q:r02-reg-applied-07}

\noindent\textbf{原典：} sources/R02/17\_applied\_science.pdf、PDF物理3頁。

\textbf{公式正答：} 4（single）。

\par\textit{原典本文の校正と独自解説は未作成。分類のみの登録項目。}

\subsection{令和3年度 問5}
\label{q:r03-reg-applied-05}

\noindent\textbf{原典：} sources/R03/17\_applied\_science.pdf、PDF物理2頁。

\textbf{公式正答：} 3（single）。

\par\textit{原典本文の校正と独自解説は未作成。分類のみの登録項目。}

\subsection{令和6年度 問5}
\label{q:r06-reg-applied-05}

\noindent\textbf{原典：} sources/R06/17\_applied\_science.pdf、PDF物理2頁。

\textbf{公式正答：} 5（single）。

\par\textit{原典本文の校正と独自解説は未作成。分類のみの登録項目。}

